\documentclass[11pt]{article}
\usepackage[margin=1in]{geometry}
\usepackage[T1]{fontenc}
\usepackage{lmodern,amsmath,amssymb,amsthm,mathtools,microtype}
\usepackage{xcolor}
\usepackage[colorlinks=true,linkcolor=blue!55!black,citecolor=blue!55!black,urlcolor=blue!55!black]{hyperref}
\hypersetup{pdftitle={Optimal Search and Positive-Adversary Bounds with Accidentally Collapsing Measurements},pdfauthor={Daniel Fredriksen},pdfsubject={Research preprint}}
\hyphenation{Milo-schew-sky}
\setlength{\emergencystretch}{1.5em}
\newtheorem{theorem}{Theorem}
\newtheorem{lemma}[theorem]{Lemma}
\newtheorem{corollary}[theorem]{Corollary}
\theoremstyle{remark}
\newtheorem{remark}[theorem]{Remark}
\newcommand{\eps}{\varepsilon}
\newcommand{\ket}[1]{\lvert #1\rangle}
\newcommand{\bra}[1]{\langle #1\rvert}
\newcommand{\ip}[2]{\langle #1\mid #2\rangle}
\newcommand{\Tr}{\operatorname{Tr}}
\title{Optimal Search and Positive-Adversary Bounds with Accidentally Collapsing Measurements}
\author{Daniel Fredriksen\\Quantyra Inc.\\\small\href{https://orcid.org/0009-0002-9123-7435}{ORCID: 0009-0002-9123-7435}}
\date{November 15, 2026\\\small Research preprint; version 0.3.0}
\begin{document}
\maketitle
\begin{abstract}
We study unstructured search on $N\ge2$ items with a hypothetical Born sampler that returns a
measurement outcome while independently projecting the sampled register with
probability $\eps$. Sampler outcomes and collapse flags are immediately
available to control operations on the same continuing state.
With a noiseless search oracle, arbitrary entangled workspace,
and ordinary partial measurements, we prove that constant-success unique-item
search has optimal total oracle-query, sampler, preparation, and ordinary
measurement cost
$\Theta(N^{1/3}+\sqrt{\eps N})$, uniformly for $0\le\eps\le1$.
Each whole-register readout counts as one call; the construction separately
uses $O(P\log N)$ transcript bits for $P$ readouts.
The upper bound uses partial Grover search followed by repeated readouts and
charges every predetermined attempt. For the lower bound, a root-fidelity
potential on complete conditional histories remains monotone through any
number of readouts when its copy exponent is $\lceil1/\eps\rceil$.
Its query sensitivity gives $Q^2\lceil1/\eps\rceil=\Omega(N)$;
a decreasing-exponent argument extends the ideal inequality
$Q^2(P+1)=\Omega(N)$ to immediate feedback. Finite-tree induction and query
frontiers handle unequal sampling and stopping histories.
Full cube-root scaling survives precisely when
$\eps=O(N^{-1/3})$. We also prove, for every finite Boolean promise problem $f$,
$Q\min\{P+1,\lceil1/\eps\rceil\}\ge\operatorname{Adv}_{+}(f)/36$,
using $P+1$ when $\eps=0$. A positive-adversary sum over complete histories
charges each synchronized query without a factor for the number of branches.
The results concern this specified postquantum access model,
not physical noisy quantum hardware or ordinary BPP versus BQP.
\end{abstract}

\section{Introduction and contribution}
Non-collapsing measurements provide classical information about a quantum state
without the usual disturbance. Aaronson, Bouland, Fitzsimons, and Lee (ABFL)
introduced the PDQP framework and a search algorithm combining partial Grover
evolution with repeated sampling \cite{abfl}. Miloschewsky and Podder proved
the optimal cube-root search lower bound and developed related
conditional-copy and fidelity methods \cite{mp}. We ask how the search cost
changes when the non-collapsing primitive sometimes performs an ordinary
projection instead.

Our primitive samples a specified register in the computational basis. An
independent Bernoulli coin determines whether that register is projected or
the global state is left intact. The sampled outcome and coin are recorded
and immediately available to control later operations. This is a particular noisy
extension of a hypothetical information resource. It does not assert a
physical implementation of non-collapsing measurement.

The noise location matters. Our oracle is noiseless, and an algorithm may
finish all its Grover queries before invoking the sampler. Noise on oracle
calls gives a different search problem \cite{rs,rosmanis}. Nor does collapse
of a sampled register necessarily destroy useful information in the other
registers. A lower bound must account for partial measurements, entanglement,
and the full classical history rather than counting apparent restarts.

\begin{theorem}[Uniform optimal cost]\label{thm:main}
In the model of Section~\ref{sec:model}, for $N\ge2$, let $T_*(N,\eps)$ be the minimum
worst-case cost for outputting the unique marked item with probability at
least $2/3$ for every marked item. Then, with absolute implied constants,
\begin{equation}\label{eq:main}
 T_*(N,\eps)=\Theta\bigl(N^{1/3}+\sqrt{\eps N}\bigr),
 \qquad 0\le\eps\le1.
\end{equation}
More precisely, for $N\ge12$, every such algorithm satisfies
\begin{equation}\label{eq:two-lower}
 Q^2(P+1)\ge\frac N{32},\qquad
 Q^2\lceil1/\eps\rceil\ge\frac N{32}\quad(\eps>0).
\end{equation}
Here $Q$ and $P$ are worst-case oracle and sampler call caps; total cost also
charges fresh attempts and ordinary measurement/reset calls. The achieving
construction uses $O(\log N)$ quantum space and
$O((N^{1/3}+\sqrt{\eps N})\log N)$ nonoracle gates and classical scheduling/scan
work, plus $\operatorname{poly}(\log N)$ parameter arithmetic,
with transcript storage charged separately. The same bounds hold if collapse
flags are hidden.
\end{theorem}

The main contribution is the fixed-exponent comparison that makes the second
inequality in~\eqref{eq:two-lower} valid despite arbitrarily many readouts.
The proof uses complete conditional histories and a mathematical copy
exponent of order $1/\eps$. The collapsing branch compensates for the
information released by the non-collapsing branch at this exponent. A Bures
distance hybrid argument then incurs only the square root of the exponent
per query. These copies are a proof device, not an additional algorithmic
resource. Since version 0.2.0, both comparisons extend to immediate feedback via
finite-tree induction and query frontiers. The original deferred-history
model is a special case; the scalar inequalities and matching upper algorithm
are inherited from version 0.1.0.
Section~\ref{sec:lower} proves the lower bounds, and
Section~\ref{sec:upper} supplies a matching fully charged construction.

Our second main result extends the same local power comparisons to
arbitrary finite Boolean promise problems through positive adversary
certificates. Define
\[
 K_\eps(P)=
 \begin{cases}
 P+1,&\eps=0,\\
 \min\{P+1,\lceil1/\eps\rceil\},&0<\eps\le1.
 \end{cases}
\]
\begin{theorem}[Positive-adversary bound]\label{thm:adversary}
Let $X\subseteq\{0,1\}^n$ be finite and $f:X\to\{0,1\}$.
In the model of Section~\ref{sec:model}, replace the search oracle by
$O_x\ket{i,b,z}=(-1)^{b x_i}\ket{i,b,z}$ and require success at least $2/3$
on every promised input. Every finite capped algorithm satisfies
\begin{equation}\label{eq:adversary-main}
 QK_\eps(P)\ge c_*\operatorname{Adv}_{+}(f)
 \ge \operatorname{Adv}_{+}(f)/36,
 \qquad c_*:=\frac{1-2\sqrt2/3}{2}>\frac1{36}.
\end{equation}
Here $\operatorname{Adv}_{+}$ is the entrywise-nonnegative spectral
adversary, defined in Section~\ref{sec:adversary}. Immediate feedback,
arbitrary entangled workspace, partial measurements, the specified resets,
and unequal stopping remain allowed.
\end{theorem}

This is a general lower-bound tool inside the declared model, with
potentially exponential analytical certificates. It makes no efficient
certificate-discovery claim. The sharper search bound in
Theorem~\ref{thm:main} is retained. Section~\ref{sec:adversary} proves the
new statement and gives a parity application.

\subsection{Relation to prior work}
ABFL supplies the partial-Grover/repeated-readout template \cite{abfl}.
Miloschewsky--Podder supplies the ideal cube-root result and closely related
copy/fidelity structure \cite{mp}. We reconstruct the ideal inequality and
extend its control scope; neither the deferred endpoint nor the use of
conditional copies is a separate contribution. The fixed-power comparison keeps
the exponent fixed as a function of $\eps$, instead of decreasing a
readout-dependent exponent. Their paper also identifies noisy computation
with non-collapsing measurements as a broader open direction.

For general Boolean promises, Miloschewsky--Podder already proves the
positive weighted bound $QP=\Omega(1/v_{\max})$ in its ideal PDQP model
\cite[Lemma 4.3]{mp}. The standard nonnegative spectral convention and
its equivalence to positive weighted methods are due to prior adversary
work, as collected by \v{S}palek--Szegedy \cite[Theorem 3.1]{ss}.
Theorem~\ref{thm:adversary} combines that machinery with our local power
comparisons and a probability-weighted query-frontier charge. The
independent-collapse and persistent-feedback scope is the extension;
neither the adversary framework nor the spectral certificate is new.
The cited ideal theorem does not already assert this noisy-feedback bound.

The author-hosted ABFL manuscript separately defines persistent adaptive
CQP: a controller receives samples and chooses further gates on the same
state \cite[Sections 3, 6, 8]{abflcqp}. Its search lower-bound question and
the feedback question in \cite[Section 1.5]{mp} concern this continuation.
At $\eps=0$, our finite capped model includes these search procedures and
answers that search question under the explicit charged-call convention.
This is search-model inclusion, not equality of entire complexity classes.
Adaptive calls to a \emph{stateless} circuit-sampling oracle in
\cite[footnote 6, Section 5]{msy} start independent invocations; their
sampling-hardness equivalence does not supply this persistent-state bound.

Abbas, King, Huang, Huggins, Movassagh, Gilboa, and McClean already showed
that sufficiently gentle imperfect measurements can support cube-root
search \cite[Appendix E.4, Proposition 28]{abbas}. Their construction uses
multiple copies with a gentleness assumption on the resulting state: the
cited proof uses $O(\log N)$ copies and $1/N$-gentleness. A
small conditional disturbance is different from a rare event with potentially
order-one conditional disturbance. Thus robustness to imperfect readout is
an existing idea; the present result gives a matched, all-$\eps$ law for the
explicit independent-collapse instrument below.

Regev--Schiff studies faulty oracle calls that may act as the identity
\cite{rs}; Rosmanis studies dephasing and depolarization associated with oracle
calls, including recorded noise events \cite{rosmanis}. Their noise acts
during acquisition of oracle information. Here the retained nondisturbing
Born record is an additional, nonstandard resource, and oracle calls remain
noiseless. Section~\ref{sec:model} explains why discarding that record changes
the operation. Fiorentino--Weigert considers a passive-measurement theory
with different composition rules and distinguishes oracle calls from
additional measurement effort \cite{fw}. Its state-reconstruction discussion
does not give the charged independent-collapse tradeoff in~\eqref{eq:main}.

Our focused comparison did not locate the same theorem or a direct
cost-preserving implication from these results. This is a bounded literature
comparison, not an assertion of exhaustive priority. Standard root fidelity,
the Bures metric, Jensen's inequality, and weighted AM--GM remain inherited
tools \cite{gln}.

\section{Access model and costs}\label{sec:model}
For the search result, the input is a unique marked label $w\in[N]=\{1,\ldots,N\}$. The controlled
phase oracle acts as
\begin{equation}
 O_w\ket{i,b,z}=(-1)^{b[i=w]}\ket{i,b,z}.
\end{equation}
A bit query is obtained with Hadamards on $b$. The oracle $O_0=I$ denotes an
all-zero reference input used only in the proof. No success requirement is
imposed on that input. Initial states, gates, and classical randomness are
independent of $w$; every oracle-dependent preparation pays its queries.

For Theorem~\ref{thm:adversary}, use instead the promised string $x\in X$
and the controlled bit-phase oracle stated there. All other operations and
initial preparations are input-independent, and the same sampler and
controller rules below apply. Padding query addresses have value zero.
There is no additional access to $x$ or free input-dependent preparation.
In that theorem the caps are maxima over promised strings and branches;
undefined formal histories are extended arbitrarily and truncated at the
same caps. No off-promise success assumption is used.

For a specified register $S$, put
$\Pi_y=\ket{y}\bra{y}_S\otimes I$. On a conditional normalized global pure
state $\psi$, a sampler call chooses $c\sim\operatorname{Bernoulli}(\eps)$
independently, draws $y$ with probability $p_y=\|\Pi_y\psi\|^2$, records
$(y,c)$, and replaces the state by
\begin{equation}\label{eq:instrument}
 \psi\quad(c=0),\qquad
 \frac{\Pi_y\psi}{\sqrt{p_y}}\quad(c=1).
\end{equation}
Zero-probability branches are omitted. The identity in $\Pi_y$ includes all
workspace and any inaccessible reference. This specifies the update on an
entangled state, not just the distribution on $S$.

A finite capped classical controller receives sampler outcomes and flags
immediately. Based on its history it may select gates, oracle queries,
sampled registers, ordinary partial projective measurements, preparations,
resets, restarts and termination, on the same continuing state. At equal
histories its decision rule is the same in all oracle worlds. Fresh classical
random choices have finite support and are input-independent given the history.
Deferred feedback
and hidden flags are restrictions of this model. There is no unknown-state
description, physical cloning, rewinding, chosen-outcome projection or free
postselection. A final classical map of the available transcript produces
the answer.

Mixed initial preparations are represented by an input-independent global
purification; classical preparation choices and ordinary measurement
outcomes are included in the history. Environments may be retained as
inaccessible purification registers. The primitive is defined on these
conditional global pure states. We do not assume an ensemble-independent
nonlinear map on arbitrary reduced density matrices.

Reset operations have specified implementations. A measurement reset performs
an ordinary projective measurement, includes its outcome in the classical
history, and prepares a known basis state. An environment reset swaps the
register with a fresh known ancilla and retains the old register as an
inaccessible quantum environment, without creating a classical outcome.
Discards similarly retain their quantum purification without conditioning
on a fictitious outcome. These implementations can produce different later
sampler statistics and are not identified as equivalent channels here.
Both reset implementations are allowed and charged as ordinary reset calls;
their gate and ancillary-space costs are additional. The lower bound covers
each implementation separately, and the main upper construction uses no resets.

To see the distinction from ordinary dephasing, write
$\rho=\ket{\psi}\bra{\psi}$ on a pure history branch. The unnormalized states
associated with the classical outcomes are
\begin{equation}\label{eq:record}
 \omega_{0,y}(\rho)=(1-\eps)\Tr(\Pi_y\rho)\rho,
 \qquad
 \omega_{1,y}(\rho)=\eps\Pi_y\rho\Pi_y.
\end{equation}
Their sum, with the record discarded, is
$(1-\eps)\rho+\eps\sum_y\Pi_y\rho\Pi_y$.
Retaining the record instead releases a Born sample while leaving the state
unchanged on the $c=0$ branch. That record-retaining rule is not an ordinary
linear quantum instrument for $\eps<1$. Equal discarded-record marginals
therefore do not justify transferring an ordinary noisy-oracle theorem.

Let $Q,P,R,M$ be the smallest worst-case caps over valid marked inputs
and complete branches on oracle calls, sampler calls, fresh attempt
preparations and ordinary projective measurement/reset calls, respectively.
Thus these caps are the corresponding branch maxima. We charge
\begin{equation}\label{eq:cost}
 T=Q+P+R+M.
\end{equation}
This sum of caps lies between the worst-case summed call count and four
times that count. The finite controller and caps apply on every execution,
including the reference oracle $O_0$. Enforce them by halting before a cap
would be exceeded; if initially promised only on marked inputs, this
truncation leaves every valid-input execution unchanged.
All actual attempts are paid, including failed trials and
feedback-selected restarts. The achieving construction schedules all attempts
in advance and pays them even after success. No expected-cost or unbounded
stopping-time theorem is claimed.
Known zero ancillas have no input information and are charged in gate and
space accounting. Register widths, transcript storage, ordinary gates,
classical processing, and precision are additional costs stated explicitly.
One wide-register readout is one call, not free printing of its output bits.
The parameter $\eps$ is known; a sufficient approximation contract is given
in Section~\ref{sec:resources}. Unknown-rate calibration is not free access.

\section{Complete-history lower bound}\label{sec:lower}
\subsection{A conditional-copy metric}
At a prefix-free cut of the finite execution tree, let $h$ include all
preceding ordinary outcomes, sampler
outcomes and flags, and classical random choices. In oracle world $a$, let
$p_a(h)$ be its probability and $\psi_{a,h}$ the conditional global pure
state. The cut includes halted leaves. Histories retain all actual classical
records, including erased records or hidden flags; they never measure an
inaccessible purification to invent a record.
Comparisons at the same $h$ apply the same input-independent
operations. Branches having zero weight in one world contribute zero.

For an integer $k\ge1$, define the mathematical state
\begin{equation}
 \Sigma_a^{(k)}=\bigoplus_h p_a(h)
 \bigl(\ket{\psi_{a,h}}\bra{\psi_{a,h}}\bigr)^{\otimes k}.
\end{equation}
Its root fidelity with $\Sigma_b^{(k)}$ is
\begin{equation}\label{eq:potential}
 F_k(a,b)=\sum_h\sqrt{p_a(h)p_b(h)}
       |\ip{\psi_{a,h}}{\psi_{b,h}}|^k.
\end{equation}
We use root fidelity rather than its square. The associated Bures distance
$d_k(a,b)=\sqrt{2-2F_k(a,b)}$ obeys the triangle inequality \cite{gln}.
For completeness, fix a purification of the middle state and choose
overlap-maximizing purifications of each endpoint. The Euclidean triangle
inequality on these purifications, followed by minimization between the
endpoints, proves the Bures triangle inequality. Conditional powers in
\eqref{eq:potential} lie \emph{inside} the history sum. No independence
of histories or of successive snapshots is assumed.

We also use a history-dependent positive integer exponent $\kappa(h)$.
Replacing $k$ by $\kappa(h)$ inside each direct-sum block defines
$\Sigma_a^{(\kappa)}$ and its root fidelity $F_\kappa(a,b)$ on the same
history-indexed block spaces in both oracle worlds. The associated Bures
distance is again a metric. This variation aligns feedback branches with
different numbers of intermediate readouts. Bures triangles are used within
one common cut and block space. Across different cuts we use scalar
monotonicity, not a metric between spaces with different copy exponents.

\begin{lemma}[Partial projection comparison]\label{lem:projection}
For normalized $u,v$ and a projective measurement $\{\Pi_y\}$, write
\[
 r=|\ip u v|,\quad p_y=\|\Pi_yu\|^2,\quad q_y=\|\Pi_yv\|^2,
 \quad r_y=\frac{|\bra u\Pi_y\ket v|}{\sqrt{p_yq_y}}.
\]
Omit terms with $p_yq_y=0$, and set
$A=\sum_y\sqrt{p_yq_y}$ and
$C_k=\sum_y\sqrt{p_yq_y}r_y^k$. For $k\ge1$ and $A>0$,
\begin{equation}\label{eq:jensen}
 r\le A\le1,\qquad C_k\ge\frac{r^k}{A^{k-1}}\ge r^k.
\end{equation}
If $A=0$ then $r=0$.
\end{lemma}
\begin{proof}
Cauchy--Schwarz gives $A\le1$ and $r_y\le1$; when $p_yq_y=0$, it also gives
$|\bra u\Pi_y\ket v|=0$, justifying omission. Completeness and the triangle
inequality give
$B_0:=\sum_y|\bra u\Pi_y\ket v|\ge r$, while $B_0\le A$.
Weighted Jensen with weights $\sqrt{p_yq_y}/A$ gives
$C_k\ge A(B_0/A)^k\ge r^k/A^{k-1}$. The zero case follows from $r\le A$.
\end{proof}

An ordinary measurement replaces the contribution $r^k$ of a common
history by $C_k$. Thus it cannot decrease $F_k$. Common unitaries preserve
$F_k$, including gates chosen after sampler outcomes and flags. Different
child histories may choose different gates; matching child histories in the
two worlds choose the same gate.
Preparations of known ancillas and environment resets are isometries on the
retained purification. Measurement resets consist of an ordinary measurement
and known preparation, with the outcome retained in the history. The two
distinct reset conventions therefore both obey the comparison. Each argument
also applies blockwise at a history-dependent exponent. Fresh classical
random branches multiply each pair contribution by common probabilities
whose sum is one, so they preserve it exactly.

\subsection{Fixed exponent for noisy readouts}
\begin{lemma}[Fixed-power monotonicity]\label{lem:fixed}
If $0<\eps\le1$ and $k$ is an integer with $\eps k\ge1$, a sampler call
cannot decrease $F_k$.
\end{lemma}
\begin{proof}
At a common history, the non-collapsing outcomes contribute
$(1-\eps)Ar^k$; the collapsing outcomes contribute $\eps C_k$.
For $A>0$, Lemma~\ref{lem:projection} gives a total of at least
\begin{equation}\label{eq:fixed}
 r^k\bigl((1-\eps)A+\eps A^{1-k}\bigr)
 \ge r^k A^{1-\eps k}\ge r^k.
\end{equation}
The first inequality is weighted AM--GM; the second uses $A\le1$ and
$\eps k\ge1$. If $A=0$ the old contribution is zero. Multiplying by the
history weight and summing proves the claim. The independent flag
probabilities produce exactly the coefficients $1-\eps$ and $\eps$.
\end{proof}

Choosing $k=\lceil1/\eps\rceil$ works for arbitrarily many readouts.
This is a potential comparison, not a simulation using $k$ physical copies.
In particular, it does not replace the actual number of readouts by their
mean waiting time to a collapse.

\subsection{Decreasing exponent for the ideal baseline}
For a separate argument, set $\kappa(h)=P+1-p(h)$, where $p(h)$ is the number
of sampler calls preceding history $h$. A real sampler call decreases that
history's exponent by one. The worst-case cap $p(h)\le P$ ensures that its
exponent is at least two immediately before a call and at least one afterwards.
Writing $k=\kappa(h)$ before the call, the new contribution is
\begin{equation}\label{eq:decreasing}
 (1-\eps)Ar^{k-1}+\eps C_{k-1}
 \ge(1-\eps)r^k+\eps r^{k-1}\ge r^k.
\end{equation}
Here $A\ge r$ and Lemma~\ref{lem:projection} apply. Ordinary measurements
retain the current exponent. If $A=0$, the old contribution is zero.
Consequently this changing potential is
nondecreasing at nonquery steps, including at $\eps=0$. This reconstructs
the ideal-search inequality with complete ordinary and sampler histories,
now allowing their immediate use for control.

\subsection{Query sensitivity}
Compare $O_0$ with each $O_w$, and abbreviate $d_w=d_k(0,w)$.
At a query, let $a_{0,h}(w)$ be the probability that the reference state
has address $w$ and oracle control bit $1$. Put
$a_0(w)=\sum_h p_0(h)a_{0,h}(w)$, so $\sum_w a_0(w)\le1$.
As a proof device, apply $O_w$ to every copy of the reference state.
Its overlap with the old reference block has magnitude
$|1-2a_{0,h}(w)|^k$. For $0\le a\le1$,
\begin{equation}\label{eq:self}
 1-|1-2a|^k\le2ka.
\end{equation}
For $a\le1/2$ this is Bernoulli's inequality; for $a\ge1/2$ the right side
is at least one. The Bures distance from the reference state to this
queried counterpart is at most $2\sqrt{k a_0(w)}$.

Applying the same $O_w$ to both mathematical comparison states preserves
their distance. The actual reference query is the identity. The triangle
inequality therefore gives
\begin{equation}
 d'_w\le d_w+2\sqrt{k a_0(w)}.
\end{equation}
By Minkowski's inequality,
\begin{equation}\label{eq:D}
 D:=\Bigl(\sum_wd_w^2\Bigr)^{1/2}
 \quad\text{increases by at most }2\sqrt{k}\text{ per query}.
\end{equation}
For history-dependent exponents, the same calculation uses
$K=\max_h\kappa(h)$: summing~\eqref{eq:self} with history weights bounds
the reference-query distance by $2\sqrt{K a_0(w)}$, so the aggregate
increment is at most $2\sqrt K$.

To handle outcome-controlled sampling, gates and stopping, align each branch
at the prefix-free frontier
immediately before their $j$th oracle call. From the previous query frontier,
advance each branch through its intervening nonquery operations to this
frontier; a halted branch is padded with identity queries using control bit
zero. The stopping location is determined by the recorded history, not by
the oracle label separately. Each intervening operation replaces one history
contribution by a sum at least as large, so their full finite subtree
expansion cannot decrease the fidelity. All branches at the frontier then
make their $j$th query together, with $\sum_w a_0(w)\le1$. This gives at most
$Q$ query frontiers, rather than summing increments over unrelated wall-clock
slots. The exponent $\kappa(h)=P+1-p(h)$ tracks each branch's actual readouts
and always satisfies $1\le\kappa(h)\le P+1$.
This is induction on the actual finite execution tree: each nonquery node
expansion obeys the local comparison, including history-selected projectors
and zero-weight histories. There is no branch-count loss. Control-zero
padding is analytical, adds no actual calls, and retains halted leaves.
The frontiers do not reschedule the algorithm. Finite controller/register
caps provide common padded quantum spaces at each cut; inaccessible
environments remain quantum.

Initially $D=0$. After the last query, advance through all remaining nonquery
operations to termination using the same monotonicity. Consequently
$D\le2Q\sqrt{\lceil1/\eps\rceil}$ in the fixed-power argument and
$D\le2Q\sqrt{P+1}$ in the decreasing-power argument.

\subsection{Success requires distance}
Under $O_0$, let $z_w$ be the probability that the final output is $w$.
Since $\sum_wz_w\le1$, at least $N-6$ labels satisfy $z_w\le1/6$.
For any such $w$, success under $O_w$ makes the two output distributions
differ by at least $1/2$ in total variation. If their classical root
affinity is $A_{\rm out}$, Cauchy--Schwarz gives
\[
 \operatorname{TV}\le\sqrt{1-A_{\rm out}^2},
 \qquad A_{\rm out}\le\sqrt{3}/2.
\]
Equation~\eqref{eq:potential} is at most the affinity of the complete
histories. Classical processing can only increase affinity, so
$F_k(0,w)\le A_{\rm out}$, also with $F_\kappa$ in place of $F_k$.
Any final ordinary quantum measurement is
included among the ordinary history steps before this processing. Hence
$d_w^2\ge2-\sqrt3>1/4$. For $N\ge12$,
\begin{equation}
 D^2\ge\frac{N-6}{4}\ge\frac N8.
\end{equation}
Combining this with~\eqref{eq:D} proves~\eqref{eq:two-lower}.
Since $\lceil1/\eps\rceil\le2/\eps$, we obtain
\[
 Q\ge\frac{\sqrt{\eps N}}8.
\]
Also $Q^2(P+1)\le2(Q+P)^3$ when $Q+P\ge1$, and therefore
$Q+P\ge N^{1/3}/4$. Preparations and ordinary measurements can only add to
the cost. The maximum and sum of the two lower bounds differ by at most a
factor of two. For $2\le N<12$, a zero-query output law is independent of
$w$ and succeeds for some $w$ with probability at most $1/N\le1/2$.
Thus $T\ge Q\ge1$, whereas $N^{1/3}+\sqrt{\eps N}<6$ in this bounded range.
This completes the uniform lower bound of Theorem~\ref{thm:main} for $N\ge2$.
The trivial $N=1$ problem is excluded, as it needs no input access.

\section{Matching search construction}\label{sec:upper}
\subsection{Exact success probability of an attempt}
Prepare a Grover state using $t$ iterations. Its marked weight is
$p=\sin^2((2t+1)\arcsin N^{-1/2})$. One additional bit query appends the
mark bit $[i=w]$. Perform $r$ sampler calls on the address and mark-bit
registers, with no intervening operation on them. A returned mark bit $1$
identifies the answer without a later outcome-dependent verification query.

Before the first collapse, the samples have the original distribution.
After a collapse, all further samples repeat that address and mark bit.
First success at draw $j+1$ requires $j$ unmarked, non-collapsing draws
and then a marked draw. Consequently the success probability is exactly
\begin{equation}\label{eq:success}
 s(p,r,\eps)=p\sum_{j=0}^{r-1}[(1-p)(1-\eps)]^j
 =\frac{p\{1-[(1-p)(1-\eps)]^r\}}{p+\eps-p\eps}.
\end{equation}
The sum defines the zero-denominator case. It yields
$1-(1-p)^r$ at $\eps=0$ and $p$ at $\eps=1$. No failed collapse is removed
by conditioning.

\subsection{Uniform parameter choice}
Pad to the least power of four $N'\ge N$, so $N'<4N$. Padded labels are
unmarked. Controlled oracle access implements this padding with an
$O(\log N)$-gate range check and one oracle call. For $N'\ge256$, set
\begin{equation}\label{eq:parameters}
 m=\left\lceil\max\{(N')^{1/3},\sqrt{\eps N'}\}\right\rceil,
 \qquad t=\lfloor m/4\rfloor,
 \qquad r=\lceil N'/m^2\rceil.
\end{equation}
Because $\sqrt{N'}$ is integral, $4\le m\le\sqrt{N'}$.
With $\theta=\arcsin(1/\sqrt{N'})$, we have
$1/\sqrt{N'}\le\theta\le2/\sqrt{N'}$,
$2t+1\ge m/4$, and
$(2t+1)\theta\le1+2/\sqrt{N'}<\pi/2$.
Using $\sin x\ge2x/\pi$ on this interval gives
\begin{equation}\label{eq:pbound}
 p\ge\frac{m^2}{4\pi^2N'}\ge\frac{m^2}{40N'},
 \qquad \eps\le40p,\qquad rp\ge1/40.
\end{equation}
In~\eqref{eq:success}, the denominator is at most $41p$, and the numerator
factor in braces is at least $1-e^{-rp}$. Thus
\begin{equation}
 s\ge\frac{1-e^{-1/40}}{41}\ge\frac1{3280}.
\end{equation}
For the last bound, $1-e^{-x}\ge x/2$ for $0\le x\le1$ suffices.
Run $B=13120$ independent fresh attempts, scheduled in advance. Their
failure probability is at most $(1-1/3280)^B\le e^{-4}<1/12$.
Independence here is justified by the explicitly fresh product states and
independent sampler coins; it was not assumed in the lower bound.
All $B$ attempts are executed and charged. For bounded $N'<256$,
exhaustively querying the bounded number of addresses is a constant-cost
fallback, using ordinary measurements and no sampler assumption.

For the main construction,
\begin{equation}\label{eq:resources}
 Q=B(t+1),\quad P=Br,\quad R=B,\quad M=0.
\end{equation}
Since $m=O(N^{1/3}+\sqrt{\eps N})$ and $N'/m^2\le(N')^{1/3}$,
these counts prove the upper half of~\eqref{eq:main}. The final classical
scan selects any recorded marked address. It uses no collapse flag, so
hiding flags changes neither the construction nor the lower bound, which
already granted access to them.

\subsection{Gate, space, and precision costs}\label{sec:resources}
All attempts can use disjoint fresh registers in one fixed circuit. As $B$
is an absolute constant, including their inactive registers takes
$O(\log N)$ quantum space. Uniform preparation over $N'$ uses Hadamards;
the Grover reflection and padding controls use $O(\log N)$ reversible gates
and clean ancillas per query block. An exact implementation using Hadamard,
Pauli $X$, and Toffoli gates exists. Nonoracle gate count and final classical
scheduling/scan time are $O((T+1)\log N)$; the full transcript requires
$O(P\log N)$ bits. The parameter arithmetic below is charged additionally.
The black-box cost of implementing a search oracle in an application must
be added separately. The constants above certify an asymptotic law and are
not a practical performance estimate.

An upper estimate $\bar\eps\in[\eps,\min(1,\eps+1/N')]$ can replace $\eps$
in~\eqref{eq:parameters}. The lower marked-weight and success guarantees
still hold, while $\sqrt{\bar\eps N'}\le\sqrt{\eps N'+1}$ changes only
constants. This requires $O(\log N)$ parameter bits, not exact arbitrary-real
arithmetic. For a supplied rational upper estimate, $m$ is the least integer satisfying
$m^3\ge N'$ and $m^2\ge\bar\eps N'$; binary search uses exact integer
comparisons and time polynomial in $\log N$. It introduces no input-dependent advice.

For approximate preparation, impose the error bound on the actual conditional
global pure state of each attempt, including its workspace and reference:
on every classical preparation branch it has trace-distance error at most
$\eta$ from the ideal preparation before the first readout. Attempts still
use independent preparations and coins as above, with approximation confined
to their individual preparation gates. The initial address/mark-bit
distributions on each such branch then have total-variation distance at most
$\eta$. Couple the collapse coins and each pre-collapse Born draw. Once a
collapse occurs, the sampled address and bit repeat, so the full transcript
of one attempt differs with probability at most $r\eta$. Across $B$
attempts the loss of success is at most $Br\eta$. Taking
$\eta\le1/(24Br)$ preserves success above $2/3$.
An error bound only on an averaged reduced density matrix does not imply this
transcript estimate for the nonlinear sampler. If $G$ within-attempt unitary
preparation gates are approximated, a sufficient condition is a sum
of operator-norm errors at most $\eta$. Per-gate accuracy $O(\eta/G)$ uses
$O(\log N)$ accuracy bits, with any gate-synthesis overhead charged
separately. The exact gate implementation already avoids that overhead.
This is a coupling for this particular no-intervening-operation construction,
not a general trace-distance contraction property of nonlinear samplers.

\subsection{A branching check and an insufficient comparator}
Take $u=\sqrt p\ket{00}+\sqrt{1-p}\ket{11}$ and
$v=\sqrt q\ket{00}+\sqrt{1-q}\ket{11}$, and write
$a=\sqrt{pq}$, $b=\sqrt{(1-p)(1-q)}$, $r=a+b$.
Sample the first qubit; halt on a collapse. On a non-collapsing outcome zero,
proceed to the next query. On outcome one, apply $X$ to the first qubit,
sample the second qubit, and then proceed. Branches have different stopping
times, registers and sample counts. Put $g_k=(1-\eps)r^{k+1}+\eps r$.
The fixed-power frontier value is
$\eps r+(1-\eps)ar^k+(1-\eps)bg_k\ge r^k$ when $\eps k\ge1$.
For the decreasing-power comparison with $P=2$, it is
$\eps r+(1-\eps)ar^2+(1-\eps)b[(1-\eps)r^2+\eps r]\ge r^3$.
At $p=1/5,q=4/5,\eps=1/2$, the unpowered one-step comparator fails:
$g_1=18/25<4/5=r$, while $g_2=82/125\ge16/25=r^2$.
The two frontier values are $412/625\ge r^2$ and $84/125\ge r^3$.
This is a symbolic check and negative control. The quantified local
inequalities and finite-tree induction prove the general claim.

\section{Positive adversaries with feedback}\label{sec:adversary}
This section proves Theorem~\ref{thm:adversary}. Constant or empty promise
problems have $\operatorname{Adv}_{+}(f)=0$ by convention and are immediate.
Otherwise let $\Gamma$ be a nonzero real symmetric entrywise-nonnegative
matrix indexed by $X$, with $\Gamma_{xy}=0$ whenever $f(x)=f(y)$. Put
\begin{equation}\label{eq:adv-definition}
 D_i[x,y]=\mathbf1[x_i\ne y_i],\qquad
 \Lambda_\Gamma=\max_i\|\Gamma\circ D_i\|,\qquad
 \operatorname{Adv}_{+}(f)=\sup_\Gamma
                \frac{\|\Gamma\|}{\Lambda_\Gamma}.
\end{equation}
The norm is the spectral operator norm and $\circ$ denotes entrywise
product. Every nonzero entry joins inputs differing in some bit, so
$\Lambda_\Gamma>0$. This is the positive spectral convention of
\cite{ss}, not the signed-weight general adversary.

Choose a nonnegative unit vector $v$ with
$v^T\Gamma v=\|\Gamma\|$. Existence does not require irreducibility:
the componentwise absolute values of a norm-extremizing real eigenvector
have Rayleigh quotient at least the absolute extremal eigenvalue, hence
exactly $\|\Gamma\|$. No division by an entry of $v$ is needed.
At any common prefix-free cut $H$ of the complete execution tree, set
\begin{equation}\label{eq:adv-potential}
 \Phi_\kappa(H)=\sum_{x,y}\Gamma_{xy}v_xv_y
  \sum_{h\in H}\sqrt{p_x(h)p_y(h)}
       |\ip{\psi_{x,h}}{\psi_{y,h}}|^{\kappa(h)}.
\end{equation}
Each execution reaches exactly one node of $H$, including halted leaves,
so $\sum_h p_x(h)=1$. Zero-probability states may be filled arbitrarily.
The exponent is a positive integer determined by the common history.
At the root the states are input-independent and
$\Phi_{\rm root}=v^T\Gamma v=\|\Gamma\|$.

\subsection{Two nonquery comparisons}
For the first argument take $\eps>0$ and the constant exponent
$k=\lceil1/\eps\rceil$. Lemma~\ref{lem:fixed} and the ordinary-operation
comparisons prove that every pair's history sum is nondecreasing through
a nonquery node. For the second, separate argument take
$\kappa(h)=P+1-p(h)$, using the actual number of preceding sampler calls.
Equation~\eqref{eq:decreasing} proves the corresponding comparison at
each sampler node, for every $0\le\eps\le1$; other nodes preserve the
exponent. In both cases entrywise positivity of $\Gamma$ preserves the
weighted comparison. Induction over each finite intervening nonquery
forest therefore makes $\Phi$ nondecreasing. Equal histories select equal
operations even with immediate feedback; no sampler outcome is deferred.
These are two distinct potentials, not a trajectory mixing their exponents.

\subsection{Spectral charge at a query frontier}
As in the search proof, advance each branch to its $j$th actual query or
to a halt. Carry halted leaves unchanged through later frontiers, with
analytical identity padding. There are at most $Q$ query layers.
At node $h$, let $\Pi_{h,i}$ project onto address $i$ and query control
$b=1$ in the history-selected query register. They are orthogonal and
$\sum_i\Pi_{h,i}\le I$; at a halted node set all these projectors to zero.
Common routing gates are included in the preceding nonquery forest.
For $u=\psi_{x,h}$ and $u'=\psi_{y,h}$,
\begin{equation}\label{eq:adv-query-pair}
 \bigl||\ip{u}{u'}|-|\ip{O_xu}{O_yu'}|\bigr|
 \le 2\sum_{i:x_i\ne y_i}\|\Pi_{h,i}u\|\|\Pi_{h,i}u'\|.
\end{equation}
Indeed $O_x^*O_y=I-2\sum_{i:x_i\ne y_i}\Pi_{h,i}$, so the triangle
and Cauchy--Schwarz inequalities apply, including complex phases and
entangled workspace. A query preserves the history probability and its
exponent. For $r,r'\in[0,1]$ and integer $k\ge1$,
$|r^k-r'^k|\le k|r-r'|$.

Let $K$ bound the exponents at the frontier, and define input-indexed vectors
\[
 a_x^{h,i}=v_x\sqrt{p_x(h)}\|\Pi_{h,i}\psi_{x,h}\|.
\]
Summing~\eqref{eq:adv-query-pair} with the nonnegative potential weights gives
\begin{align}
 \Phi(H_j)-\Phi(H_j^+)
 &\le 2K\sum_{h,i}(a^{h,i})^T(\Gamma\circ D_i)a^{h,i}\notag\\
 &\le 2K\Lambda_\Gamma\sum_{h,i}\|a^{h,i}\|_2^2
 \le 2K\Lambda_\Gamma.\label{eq:adv-query-charge}
\end{align}
The last step uses the complete-history normalization
\begin{equation}\label{eq:adv-history-normalization}
 \sum_{h,i}\|a^{h,i}\|_2^2
 =\sum_xv_x^2\sum_h p_x(h)
               \sum_i\|\Pi_{h,i}\psi_{x,h}\|^2
 \le\sum_xv_x^2\sum_h p_x(h)=1.
\end{equation}
Thus the number of histories creates no multiplicative loss. In particular,
different registers, unequal sampling counts and halted branches are
included with their actual probabilities. Using $K=P+1$ for varying
exponents is valid because all summands in the pairwise bound are
nonnegative. This reasoning does not justify a signed-weight extension.

\subsection{Terminal separation and the two bounds}
Include the actual final Boolean output in the terminal history. Opposite
outputs of $f$ require output distributions of total variation at least
$1/3$. Their affinity is consequently at most $\beta=2\sqrt2/3$.
For either positive exponent assignment,
\[
 F_\kappa(x,y)\le\sum_h\sqrt{p_x(h)p_y(h)}
 \le\sum_{b=0,1}\sqrt{\Pr_x[b]\Pr_y[b]}\le\beta
 \quad(f(x)\ne f(y)).
\]
The middle inequality groups histories by output and applies
Cauchy--Schwarz. The last uses the same affinity/variation inequality as
the search proof. Hence $\Phi_{\rm final}\le\beta\|\Gamma\|$.
Combining nonquery monotonicity with~\eqref{eq:adv-query-charge} yields
\[
 \beta\|\Gamma\|\ge\Phi_{\rm final}
 \ge\|\Gamma\|-2QK\Lambda_\Gamma,
 \qquad QK\ge c_*\frac{\|\Gamma\|}{\Lambda_\Gamma}.
\]
First apply this with constant $K=\lceil1/\eps\rceil$ for $\eps>0$.
Independently apply it with $K=P+1$ for all $\eps$. Optimizing over
$\Gamma$ and taking the minimum of these two valid bounds proves
Theorem~\ref{thm:adversary}. Since $2\sqrt2/3<17/18$, we have $c_*>1/36$.
At $P=0$ or $\eps=1$ the result gives an ordinary positive-adversary lower
bound. At $\eps=0$ only the decreasing-power argument is needed.

\subsection{Costs, application and a stopping check}
Let $W$ be the worst-case sum of actual query, sampler, preparation and
ordinary measurement/reset calls on a valid branch. With smallest valid
caps, $Q,P,R,M\le W\le T\le4W$. A nonconstant $f$ has
$\operatorname{Adv}_{+}(f)\ge1$ by a single opposite-output pair, and
requires $Q\ge1$. Thus $Q(P+1)\le W(W+1)\le2W^2$ and
$\lceil1/\eps\rceil\le2/\eps$ give separate square-root and
linear lower bounds. Their maximum is at least half their sum, giving
\begin{equation}\label{eq:adv-total}
 W=\Omega\!\left(\sqrt{\operatorname{Adv}_{+}(f)}
                  +\eps\operatorname{Adv}_{+}(f)\right).
\end{equation}
Every failed trial and query-dependent preparation is included. Gate,
classical, output-bit, space and precision costs remain additional; this
is not a matching total-bit-runtime assertion. Matrices, eigenvectors and
complete histories may be exponentially large analytical certificates.
The controller is not given their amplitudes or an efficient way to find them.

\begin{corollary}[Parity at a fixed positive collapse rate]
For $n$-bit parity and any fixed $\eps>0$, success at least $2/3$ requires
$\Omega(n)$ input queries, irrespective of the sampler cap. Reading every
bit gives a matching $O(n)$ query upper bound.
\end{corollary}
\begin{proof}
Use the hypercube adjacency matrix as $\Gamma$. Adjacent inputs have
opposite parity, $\|\Gamma\|=n$, and $\Gamma\circ D_i$ is a perfect
matching matrix of norm one. The fixed-power bound gives explicitly
$Q\ge n/(36\lceil1/\eps\rceil)$. The ordinary
upper algorithm pays its queries, measurements and parity computation.
\end{proof}
At $\eps=0$, the general theorem instead gives the tradeoff
$Q(P+1)\ge n/36$. This is an application of a standard certificate.
For search, the stronger
bound of Theorem~\ref{thm:main} remains the appropriate guarantee.

For a symbolic stopping check take $f(x)=x$ on one bit,
$\Gamma=\left(\begin{smallmatrix}0&1\\1&0\end{smallmatrix}\right)$ and
$v=(1,1)/\sqrt2$. Sample an input-independent $\ket{++}$ register once.
On label $00$, halt with output zero. Otherwise query address one on a
fresh response $\ket+$ and measure it in the $X$ basis to recover $x$.
The flag may be ignored by the controller but remains in the history.
Success is one on $x=0$ and $3/4$ on $x=1$, with $Q=P=1$.
After the sampler the decreasing exponent is one. The query reduces
$\Phi$ from one to $1/4$, while the query mass in
\eqref{eq:adv-history-normalization} is $(3/4)(1/2)=3/8$.
Using the actual exponent one, the charge is exactly $2(3/8)=3/4$.
Appending $t$ unused fair classical bits divides each history's probability
by $2^t$ without changing the total charge. Those bits and their work/space
are paid. This symbolic family checks the stopping bookkeeping;
the arbitrary-tree proof, not small-case agreement, establishes the theorem.

The lower bound applies to actual finite-precision input-independent gates
and preparations whenever the implemented algorithm meets the success
contract. An ideal-to-implemented accuracy transfer needs its own paid
whole-history argument; a linear-channel error hybrid is not assumed for
this nonlinear sampler. If such an argument only gives failure
$\eta<1/2$, the same proof replaces $c_*$ by
$(1-2\sqrt{\eta(1-\eta)})/2$. No efficient certificate discovery, signed
adversary theorem, physical realization or general classical simulator
follows from this result.

\section{Consequences and limitations}
\begin{corollary}[Persistent adaptive ideal search]
At $\eps=0$, immediate-feedback search in the declared finite capped model,
including the persistent CQP search procedures of \cite{abflcqp}, has
optimal charged-call cost $\Theta(N^{1/3})$.
\end{corollary}
\begin{proof}
Set $\eps=0$ in Theorem~\ref{thm:main}. The cited definition evolves a pure
state by unitary $U_i$ and ordinary partial measurement $M_i$, then samples
the resulting full-register Born distribution without changing the state.
The controller may select a later segment after receiving earlier samples.
Implement each $U_i,M_i$ directly, followed by one sampler with $c=0$;
retain actual ordinary outcomes analytically even when not reported.
Each returned sample is charged separately, not a free transcript for one
circuit-oracle call. This preserves search-query and sample counts and
adds at most one ordinary-measurement charge per sampled stage, without
copies or replay. The fixed-policy construction belongs to the subclass;
unused samples after its preparation gates, if required by that syntax,
add only its separately stated gate overhead. Thus the exact $\Theta$ law
uses the charged-call convention here; elementary-gate runtime carries the
stated logarithmic and parameter-processing overhead. No bound on oracle
queries alone with unlimited free samples is asserted.
\end{proof}

\begin{corollary}
The cost is $\Theta(N^{1/3})$ exactly for rates
$\eps=O(N^{-1/3})$. Any fixed positive rate gives $\Theta(\sqrt N)$ cost.
If $\eps=N^{-\alpha}$ for $\alpha\ge0$, the cost exponent is
$\max\{1/3,(1-\alpha)/2\}$.
\end{corollary}
\begin{proof}
Compare the two terms in~\eqref{eq:main}.
\end{proof}

\begin{remark}[Collapse does not imply global reset]
Let $D$ contain a query-dependent state and let the sampled register be a
separate $\ket0_S$. Every collapse of $S$ leaves $D$ intact. This refutes
a global-reset premise even within the declared model; it is not itself a
faster search algorithm. The complete-history potential is what makes the
lower bound survive this example and arbitrary entangled workspace.
\end{remark}

The theorem charges a particular postquantum sampler and a noiseless oracle.
It does not classify all noisy PDQP or adaptive CQP computations, establish
a physical realization, or improve general classical simulation of
ordinary quantum computation. Additional oracle, storage, or gate noise and
unknown-rate estimation would change the access contract. The interpolation
is an oracle-model result, with the ordinary BPP/BQP question left open.

\paragraph{Evidence and assistance disclosure.}
This preprint has an informal analytic proof and internal checks, not formal
certification or independent peer review. The accompanying archive preserves
two historical finite exact diagnostic runs and their source hashes; these illustrate
algebraic and projector comparisons and do not establish universal claims
by finite agreement. Those runs predate the feedback and general-adversary
extensions; both are proved analytically, without a new numerical campaign.
OpenAI Codex assisted with derivation, literature
comparison, drafting, and internal checking. Separate AI review was conducted
through provider CLIs. The archive records actual model identities, review
scopes, findings, and dispositions. These reviews are not independent human
peer review or formal certification.

\paragraph{Availability.}
LaTeX, the PDF, reproducibility records, and citation metadata are available at
\url{https://github.com/Quantyra/accidental-collapse-search}. Release and DOI details are recorded in the repository. Original archive materials are licensed
under CC BY 4.0. Cited works retain their own licenses.

\begin{thebibliography}{10}
\small
\raggedright
\bibitem{abfl}
Scott Aaronson, Adam Bouland, Joseph Fitzsimons, and Mitchell Lee.
The space ``just above'' BQP.
\emph{ITCS}, 2016. Preprint 2014.
\href{https://arxiv.org/abs/1412.6507v1}{arXiv:1412.6507v1}.

\bibitem{abflcqp}
Scott Aaronson, Adam Bouland, Joseph Fitzsimons, and Mitchell Lee.
The space ``just above'' BQP. Author-hosted CQP/naCQP manuscript,
\href{https://theory.stanford.edu/~abouland/Papers/cqp.pdf}{28-page PDF},
Sections 3, 6 and 8; inspected November 15, 2026. This undated source is
distinguished from the pinned arXiv PDQP version above.

\bibitem{msy}
Tomoyuki Morimae, Yuki Shirakawa, and Takashi Yamakawa.
Quantum cryptography and hardness of non-collapsing measurements.
\href{https://arxiv.org/abs/2510.04448v1}{arXiv:2510.04448v1}, October 6, 2025.

\bibitem{mp}
David Miloschewsky and Supartha Podder.
New lower-bounds for quantum computation with non-collapsing measurements.
\emph{CCC}, 2025.
\href{https://doi.org/10.4230/LIPIcs.CCC.2025.12}{doi:10.4230/LIPIcs.CCC.2025.12}.
The comparison uses \href{https://arxiv.org/abs/2411.04085v3}{arXiv:2411.04085v3},
September 2025.

\bibitem{ss}
Robert \v{S}palek and Mario Szegedy.
All quantum adversary methods are equivalent.
\emph{Theory of Computing} 2(1):1--18, 2006.
\href{https://doi.org/10.4086/toc.2006.v002a001}{doi:10.4086/toc.2006.v002a001}.

\bibitem{abbas}
Amira Abbas, Robbie King, Hsin-Yuan Huang, William J. Huggins,
Ramis Movassagh, Dar Gilboa, and Jarrod R. McClean.
On quantum backpropagation, information reuse, and cheating measurement collapse.
\emph{Advances in Neural Information Processing Systems} 36, 2023.
\href{https://proceedings.neurips.cc/paper_files/paper/2023/file/8c3caae2f725c8e2a55ecd600563d172-Paper-Conference.pdf}{Proceedings paper},
Appendix E.4, Proposition 28.

\bibitem{rs}
Oded Regev and Liron Schiff.
Impossibility of a quantum speed-up with a faulty oracle.
\emph{ICALP}, 2008.
\href{https://arxiv.org/abs/1202.1027v1}{arXiv:1202.1027v1} (2012 upload).

\bibitem{rosmanis}
Ansis Rosmanis.
Quantum search with noisy oracle.
\href{https://arxiv.org/abs/2309.14944v1}{arXiv:2309.14944v1}, 2023.

\bibitem{fw}
Vincenzo Fiorentino and Stefan Weigert.
A quantum theory with non-collapsing measurements.
\emph{Physics Letters A} 559, 130903, 2025.
\href{https://arxiv.org/abs/2303.13411v2}{arXiv:2303.13411v2}.

\bibitem{gln}
Alexei Gilchrist, Nathan K. Langford, and Michael A. Nielsen.
Distance measures to compare real and ideal quantum processes.
\emph{Physical Review A} 71, 062310, 2005.
\href{https://doi.org/10.1103/PhysRevA.71.062310}{doi:10.1103/PhysRevA.71.062310}.
\href{https://arxiv.org/abs/quant-ph/0408063v2}{arXiv:quant-ph/0408063v2}.
\end{thebibliography}
\end{document}
